\documentclass[12pt,a4paper]{article}

% ── Packages ────────────────────────────────────────────────────────────────
\usepackage[margin=2.5cm]{geometry}
\usepackage{hyperref}
\usepackage{booktabs}
\usepackage{array}
\usepackage{longtable}
\usepackage{xcolor}
\usepackage{mdframed}
\usepackage{enumitem}
\usepackage{titlesec}
\usepackage{fancyhdr}
\usepackage{graphicx}
\usepackage{microtype}
\usepackage{parskip}
\usepackage{tabularx}
\usepackage{colortbl}
\usepackage{amsmath}
\usepackage{fontawesome5}
\usepackage[T1]{fontenc}
\usepackage{lmodern}

% ── Colors ──────────────────────────────────────────────────────────────────
\definecolor{primary}{RGB}{45, 95, 70}
\definecolor{accent}{RGB}{109, 139, 116}
\definecolor{warning}{RGB}{200, 80, 60}
\definecolor{info}{RGB}{30, 90, 160}
\definecolor{lightgray}{RGB}{245, 245, 245}
\definecolor{darkgray}{RGB}{60, 60, 60}
\definecolor{clientbg}{RGB}{240, 248, 255}
\definecolor{hwbg}{RGB}{245, 250, 245}
\definecolor{rowalt}{RGB}{248, 248, 248}

% ── Section formatting ───────────────────────────────────────────────────────
\titleformat{\section}{\large\bfseries\color{primary}}{}{0em}{}[\color{accent}\titlerule]
\titleformat{\subsection}{\normalsize\bfseries\color{darkgray}}{}{0em}{}
\titleformat{\subsubsection}{\normalsize\itshape\color{darkgray}}{}{0em}{}

% ── Header / Footer ──────────────────────────────────────────────────────────
\pagestyle{fancy}
\fancyhf{}
\rhead{\textcolor{accent}{\small AI Toy Companion — Latency Analysis}}
\lhead{\textcolor{accent}{\small Confidential}}
\rfoot{\textcolor{darkgray}{\small Page \thepage}}
\lfoot{\textcolor{darkgray}{\small \today}}
\renewcommand{\headrulewidth}{0.4pt}
\renewcommand{\footrulewidth}{0.4pt}

% ── Custom environments ──────────────────────────────────────────────────────
\newmdenv[
  backgroundcolor=clientbg,
  linecolor=info,
  linewidth=1.5pt,
  leftmargin=0pt,
  rightmargin=0pt,
  innerleftmargin=12pt,
  innerrightmargin=12pt,
  innertopmargin=10pt,
  innerbottommargin=10pt,
  roundcorner=4pt
]{clientbox}

\newmdenv[
  backgroundcolor=hwbg,
  linecolor=primary,
  linewidth=1.5pt,
  leftmargin=0pt,
  rightmargin=0pt,
  innerleftmargin=12pt,
  innerrightmargin=12pt,
  innertopmargin=10pt,
  innerbottommargin=10pt,
  roundcorner=4pt
]{hwbox}

\newmdenv[
  backgroundcolor=lightgray,
  linecolor=warning,
  linewidth=1pt,
  leftmargin=0pt,
  rightmargin=0pt,
  innerleftmargin=12pt,
  innerrightmargin=12pt,
  innertopmargin=8pt,
  innerbottommargin=8pt,
  roundcorner=4pt
]{constraintbox}

% ════════════════════════════════════════════════════════════════════════════
\begin{document}

% ── Title Page ───────────────────────────────────────────────────────────────
\begin{titlepage}
  \centering
  \vspace*{2cm}
  {\color{primary}\rule{\linewidth}{1.5pt}}\\[0.4cm]
  {\Huge\bfseries\color{primary} AI Toy Companion}\\[0.3cm]
  {\LARGE\color{darkgray} Latency Analysis \& Architecture Options}\\[0.3cm]
  {\color{primary}\rule{\linewidth}{1.5pt}}\\[1.5cm]

  \begin{tabular}{rl}
    \textbf{Document Type:}   & Technical Report \\[4pt]
    \textbf{Status:}          & Draft for Review \\[4pt]
    \textbf{Date:}            & \today \\[4pt]
    \textbf{Prepared by:}     & Software \& Hardware Teams \\[4pt]
    \textbf{Requested by:}    & Kiran (Client) \\
  \end{tabular}

  \vspace{2cm}
  \begin{mdframed}[backgroundcolor=lightgray, linewidth=0pt, roundcorner=6pt,
    innerleftmargin=16pt, innerrightmargin=16pt, innertopmargin=14pt, innerbottommargin=14pt]
    \centering
    \textbf{Executive Summary}\\[6pt]
    \normalsize
    This document responds to the client's request for a progressive chunked transfer architecture
    to reduce end-to-end voice interaction latency. It covers the current baseline, all viable
    options proposed by both the hardware and software teams, an honest latency estimate for each,
    and a clear recommendation for the path forward.
  \end{mdframed}

  \vfill
  {\small\color{darkgray} Confidential — Internal Use Only}
\end{titlepage}

% ── Table of Contents ────────────────────────────────────────────────────────
\tableofcontents
\newpage

% ════════════════════════════════════════════════════════════════════════════
\section{Client's Request}

The following message was received from the client requesting an architecture revision:

\begin{clientbox}
\textbf{From:} Kiran \hfill \textbf{To:} Development Team\\[6pt]
Hi Izza,

Thanks for the update. A total latency of $\sim$37 seconds for a 30-second audio clip is not
acceptable for the intended user experience.

The issue appears to be that the current implementation waits for the full audio file to be
recorded and transferred before processing begins. For v1, we should not use a full-file
store-and-transfer approach.

Please revise the flow toward:

\begin{enumerate}[leftmargin=1.5em]
  \item Progressive chunk transfer during recording
  \item Start STT while audio is still being received
  \item Optimise for shorter interactive utterances (e.g.\ 5--15 seconds) rather than long
        full-file uploads
  \item Reduce post-speech delay as much as possible before PCB is finalised
\end{enumerate}

Please let me know the expected latency if chunked transfer + early STT is implemented, and
whether this can be tested on the current hardware before PCB order.

I have mentioned this as a concern from the start and have been asking for updates for quite
some time on this point — I also mentioned multiple times we need to do progressive chunking
and that sending the full file would not work.

\medskip
\textit{One additional point: if chunked transfer + early STT still doesn't bring latency down
enough, we should then evaluate a streaming / real-time STT approach as the fallback option.
But first, please optimise the current architecture around progressive chunking and early
processing so we can compare expected latency which is what I have been mentioning from the
start. Anything greater than that is not ready for handover.}

\medskip
Cheers,\\
\textbf{Kiran}
\end{clientbox}

% ════════════════════════════════════════════════════════════════════════════
\section{Current Architecture \& Baseline Latency}

\subsection{Current Data Flow}

The existing implementation uses a \textbf{store-and-forward} model:

\begin{enumerate}
  \item User presses and holds the physical button on the ESP32 hardware
  \item ESP32 reads I2S microphone data and writes it to LittleFS flash storage
  \item When the button is released, the ESP32 sends a \texttt{NEW:filename:size} notification
        over BLE to the mobile app
  \item The app responds with \texttt{SEND:filename}
  \item The ESP32 reads the file from flash and streams it in chunks over BLE
  \item The app assembles all chunks into a complete WAV file
  \item The complete WAV is sent to the Google STT Supabase edge function (batch request)
  \item The transcript is sent to the Gemini LLM Supabase edge function
  \item The LLM response is sent to the Resemble AI TTS Supabase edge function
  \item The resulting WAV audio is sent back to the ESP32 over BLE in 500-byte chunks
  \item The ESP32 plays the audio through the speaker
\end{enumerate}

\subsection{Baseline Latency Breakdown}

The following table shows the time taken at each stage for a representative 5-second utterance
recorded at 16\,kHz, 16-bit mono (160\,KB of audio data):

\begin{table}[h!]
\centering
\renewcommand{\arraystretch}{1.4}
\begin{tabularx}{\linewidth}{|l|X|r|}
\hline
\rowcolor{primary}
\textcolor{white}{\textbf{Stage}} &
\textcolor{white}{\textbf{Description}} &
\textcolor{white}{\textbf{Duration}} \\
\hline
Recording          & User holds button, ESP32 captures audio via I2S & 5.0\,s \\
\hline
\rowcolor{rowalt}
Flash write        & Firmware flushes audio buffer to LittleFS        & $\sim$0.5\,s \\
\hline
BLE transfer       & Complete WAV file streamed from ESP32 to app (ACK per chunk, 512 MTU, $\sim$30--50\,KB/s effective throughput) & $\sim$3--5\,s \\
\hline
\rowcolor{rowalt}
STT (cloud batch)  & Google STT edge function processes full audio file & $\sim$1--2\,s \\
\hline
LLM                & Gemini Pro edge function generates response        & $\sim$1--3\,s \\
\hline
\rowcolor{rowalt}
TTS                & Resemble AI edge function synthesises audio        & $\sim$1--3\,s \\
\hline
BLE return         & TTS WAV streamed from app to ESP32                 & $\sim$1--3\,s \\
\hline
\rowcolor{primary}
\textcolor{white}{\textbf{Total}} &
\textcolor{white}{\textbf{5-second utterance}} &
\textcolor{white}{\textbf{$\sim$12--21\,s}} \\
\hline
\end{tabularx}
\caption{Baseline latency for a 5-second utterance (current architecture)}
\end{table}

\begin{constraintbox}
\textbf{Why 30-second clips show $\sim$37 seconds of latency:} For a 30-second utterance the
audio data is $\sim$960\,KB. The BLE transfer of this file alone takes 18--30 seconds at current
throughput, consistent with the 37-second figure reported by the client.
\end{constraintbox}

% ════════════════════════════════════════════════════════════════════════════
\section{Hardware Team Response}

The following options were proposed by the hardware engineering team:

\begin{hwbox}
\textbf{Hardware Team — Proposed Options}\\[8pt]

\textbf{Current Baseline}\\
The current store-and-forward architecture records the full audio clip, saves it to flash, then
transfers it over BLE. For a 5-second utterance this results in approximately 12--15 seconds of
total latency before a response is heard — which we agree is not acceptable for the intended
experience.

\bigskip
\textbf{Option 1 — Aggressive Chunked Streaming (Recommended)}\\
This is the approach you have been requesting and it is now implemented in the firmware. As soon
as the button is pressed, audio chunks stream live over BLE to the app in real time — no waiting
for the recording to finish. The app feeds these chunks directly into a streaming STT API (such
as Google Speech-to-Text streaming or Deepgram) which begins transcribing while the user is still
speaking. By the time the user releases the button, the transcript is already mostly ready. The
LLM and TTS response then follows.

\textit{Pros:} Lowest latency of all options. No hardware changes required. Can be tested on
current hardware immediately.\\
\textit{Cons:} Requires the app to implement a streaming STT API rather than a batch API.
Slightly more complex app-side logic.

\bigskip
\textbf{Option 2 — On-Device LLM on the App (Phone-Side)}\\
A lightweight local LLM such as Llama~3.2 (1B or 3B parameter model) can be run directly on the
phone using frameworks like MediaPipe LLM Inference or llama.cpp for Android. This eliminates
the cloud API round-trip for the LLM step entirely, saving approximately 1--2 seconds.

\textit{Pros:} No internet required for the LLM step. Removes cloud API costs. Consistent latency
regardless of network conditions.\\
\textit{Cons:} App size increases significantly — approximately 2--4\,GB depending on model.
Requires a mid-range or higher phone (8\,GB RAM recommended). STT and TTS would still require
either cloud APIs or additional on-device models.

\bigskip
\textbf{Option 3 — Local Processing Entirely on MCU}\\
This approach moves STT, LLM, and TTS processing back onto the ESP32 or a more capable companion
MCU such as a Raspberry Pi or a dedicated edge AI module. A very small quantised model (under
500\,MB) runs locally with no phone or cloud dependency whatsoever.

\textit{Pros:} Fully offline. No app required beyond BLE control. No cloud API costs. Most
private solution.\\
\textit{Cons:} Constrained by MCU compute power — model quality and response intelligence will
be noticeably limited. 4--6 second latency is the realistic floor given current edge hardware.
Requires hardware changes before PCB is finalised.

\bigskip
\textbf{Option 4 — Hybrid Pipeline with Response Pre-fetching}\\
This option combines Option 1 with an app-side optimisation: the app begins constructing the
LLM prompt and pre-warming the TTS engine while STT is still streaming. As soon as the first
sentence of the LLM response is generated, TTS converts that sentence and begins sending it to
the ESP32 over BLE — without waiting for the full response to be generated. The user hears the
beginning of the answer within 1.5--2 seconds of finishing their utterance.

\textit{Pros:} Best perceived latency of all options — user hears a response almost immediately.
Works with existing hardware. No MCU changes required.\\
\textit{Cons:} Most complex app-side implementation. Requires careful pipeline coordination
between STT, LLM, and TTS streams.
\end{hwbox}

% ════════════════════════════════════════════════════════════════════════════
\section{Software Team Analysis}

\subsection{Key Constraints}

\begin{constraintbox}
\textbf{Non-negotiable constraints for the current implementation:}
\begin{itemize}[leftmargin=1.5em]
  \item All AI processing routes through \textbf{Supabase Edge Functions} (Deno runtime,
        stateless HTTP request--response only — no persistent WebSocket connections)
  \item STT must use either \textbf{Whisper.cpp} or \textbf{Google Speech-to-Text}
  \item Whisper.cpp is a C++ binary — it \textbf{cannot run} inside a Supabase edge function
  \item Google STT (batch API) \textbf{returns no transcription results} for audio shorter
        than $\sim$0.5 seconds — batch chunking of live audio does not work reliably
  \item The Supabase architecture \textbf{cannot support streaming STT} (WebSocket or gRPC)
        without adding a separate server
\end{itemize}
\end{constraintbox}

\subsection{Updated Option Analysis}

\subsubsection{Option 1 — Aggressive Chunked BLE Streaming}

\textbf{What changes (firmware):} Remove LittleFS file storage from the recording loop.
While the user holds the button, I2S audio data is sent directly over BLE notifications
in real time instead of being written to flash. A start marker is sent at the beginning,
raw PCM chunks follow, and an end marker is sent when the button is released.

\textbf{What changes (app):} The app accumulates BLE chunks as they arrive. The moment the
end marker is received, the \emph{complete} audio is passed to the STT pipeline (Google STT
batch). There is no post-recording BLE transfer delay because transfer happened concurrently
with recording.

\textbf{Important caveat on ``early STT'' during streaming:} The hardware team's proposal to
feed live chunks into a streaming STT API cannot be implemented within our current constraints.
Google STT batch API requires a minimum audio length to return results, and sending sub-second
fragments produces empty responses. Streaming STT (Google STT Streaming API, Deepgram) requires
a persistent WebSocket or gRPC connection which Supabase edge functions do not support.

\begin{table}[h!]
\centering
\renewcommand{\arraystretch}{1.3}
\begin{tabularx}{\linewidth}{|l|X|}
\hline
\rowcolor{primary}
\textcolor{white}{\textbf{Aspect}} & \textcolor{white}{\textbf{Detail}} \\
\hline
Latency saved        & Eliminates the $\sim$3--5\,s post-recording BLE transfer \\
\hline
\rowcolor{rowalt}
New total estimate   & $\sim$8--13\,s for a 5-second utterance \\
\hline
Feasibility          & High — firmware change only, no infrastructure change \\
\hline
\rowcolor{rowalt}
Risk                 & BLE must sustain 32\,KB/s (16\,kHz $\times$ 16-bit $\times$ 1ch) in real time; recording is lost if BLE drops during capture \\
\hline
Hardware change      & Yes — firmware refactor of recording loop \\
\hline
\rowcolor{rowalt}
App change           & Minor — accumulate chunks then run pipeline on END signal \\
\hline
\end{tabularx}
\caption{Option 1 summary}
\end{table}

\subsubsection{Option 2 — On-Device LLM on Phone}

The hardware team proposes running Llama~3.2 (1B--3B parameters) locally on the phone to
eliminate the LLM cloud round-trip ($\sim$1--2\,s saved).

\begin{table}[h!]
\centering
\renewcommand{\arraystretch}{1.3}
\begin{tabularx}{\linewidth}{|l|X|}
\hline
\rowcolor{primary}
\textcolor{white}{\textbf{Aspect}} & \textcolor{white}{\textbf{Detail}} \\
\hline
Latency saved        & $\sim$1--2\,s (cloud LLM round-trip eliminated) \\
\hline
\rowcolor{rowalt}
Local inference time & 1--3\,s (1B model), 2--5\,s (3B model) on mid-range CPU — net gain is small \\
\hline
App size impact      & +1\,GB (1B model) to +4\,GB (3B model) — unsuitable for most app stores \\
\hline
\rowcolor{rowalt}
RAM requirement      & 6--8\,GB phone RAM recommended; crashes on lower-end devices \\
\hline
Response quality     & Significantly below Gemini Pro \\
\hline
\rowcolor{rowalt}
Feasibility          & Very Low — app size alone disqualifies for consumer deployment \\
\hline
Recommendation       & \textbf{Do not pursue for v1} \\
\hline
\end{tabularx}
\caption{Option 2 summary}
\end{table}

\subsubsection{Option 3 — Local Processing on MCU}

Moving STT, LLM, and TTS entirely onto an edge device (Raspberry Pi, ESP32-S3, or dedicated AI
module) eliminates all cloud and phone dependencies but at a significant quality cost.

\begin{table}[h!]
\centering
\renewcommand{\arraystretch}{1.3}
\begin{tabularx}{\linewidth}{|l|X|}
\hline
\rowcolor{primary}
\textcolor{white}{\textbf{Aspect}} & \textcolor{white}{\textbf{Detail}} \\
\hline
Latency estimate     & 4--8\,s (realistic floor for edge hardware with quantised models) \\
\hline
\rowcolor{rowalt}
Model quality        & Severely constrained — tiny quantised models produce noticeably inferior responses \\
\hline
Hardware change      & Yes — significant; new MCU or companion board before PCB finalisation \\
\hline
\rowcolor{rowalt}
Offline capability   & Fully offline — no internet required \\
\hline
Feasibility          & Medium — hardware redesign required but decision must be made before PCB finalisation \\
\hline
\rowcolor{rowalt}
Recommendation       & \textbf{Evaluate in Phase 1 in parallel with Option 1} — PCB decision point \\
\hline
\end{tabularx}
\caption{Option 3 summary}
\end{table}

\subsubsection{Option 4 — Hybrid Pipeline: Chunked BLE + Sentence-Level Overlapped TTS}

This is the most impactful option the software team can implement on the \emph{app side} without
any infrastructure changes. It combines Option~1 with sentence-level pipelining:

\begin{enumerate}
  \item BLE streaming during recording (Option 1) — eliminates transfer delay
  \item When STT returns, immediately send transcript to LLM with token streaming enabled
  \item Detect the first complete sentence in the LLM token stream
  \item Send that sentence to TTS and begin streaming the resulting audio to ESP32
  \item Continue for each subsequent sentence while LLM is still generating
\end{enumerate}

The user hears the \emph{first sentence} of the response approximately 1.5--3 seconds after
releasing the button, even though the full response has not yet been generated.

\begin{table}[h!]
\centering
\renewcommand{\arraystretch}{1.3}
\begin{tabularx}{\linewidth}{|l|X|}
\hline
\rowcolor{primary}
\textcolor{white}{\textbf{Aspect}} & \textcolor{white}{\textbf{Detail}} \\
\hline
Perceived latency    & \textbf{1.5--3\,s} after button release to first audio heard \\
\hline
\rowcolor{rowalt}
Total response time  & 6--10\,s (full response, same as Option 1 but perceived start is much earlier) \\
\hline
Complexity           & High — requires LLM streaming (SSE), sentence boundary detection, concurrent TTS + BLE \\
\hline
\rowcolor{rowalt}
Feasibility          & Medium — Gemini supports streaming; Supabase edge functions support SSE; firmware must handle concurrent audio receive + play \\
\hline
Hardware change      & No (beyond Option 1 firmware change) \\
\hline
\rowcolor{rowalt}
Recommendation       & \textbf{Implement after Option 1 is stable} \\
\hline
\end{tabularx}
\caption{Option 4 summary}
\end{table}

% ════════════════════════════════════════════════════════════════════════════
\section{Latency Comparison Matrix}

The following table shows realistic end-to-end latency estimates for each option applied to a
\textbf{5-second utterance}. All times are in seconds.

\begin{table}[h!]
\centering
\renewcommand{\arraystretch}{1.5}
\begin{tabularx}{\linewidth}{|X|r|r|r|r|r|r|r|}
\hline
\rowcolor{primary}
\textcolor{white}{\textbf{Scenario}} &
\textcolor{white}{\textbf{Rec.}} &
\textcolor{white}{\textbf{Transfer}} &
\textcolor{white}{\textbf{STT}} &
\textcolor{white}{\textbf{LLM}} &
\textcolor{white}{\textbf{TTS}} &
\textcolor{white}{\textbf{BLE back}} &
\textcolor{white}{\textbf{Total}} \\
\hline
Baseline (current) & 5 & 4 & 1.5 & 2 & 2 & 2 & \textbf{16.5} \\
\hline
\rowcolor{rowalt}
Option 1: BLE streaming & 5 & 0 & 1.5 & 2 & 2 & 2 & \textbf{12.5} \\
\hline
Option 2: On-device LLM & 5 & 4 & 1.5 & 2 & 2 & 2 & \textbf{16.5}$^*$ \\
\hline
\rowcolor{rowalt}
Option 3: MCU local & 5 & 0 & 1 & 2 & 1 & 0 & \textbf{$\sim$9}$^\dagger$ \\
\hline
Option 4: Option 1 + overlapped TTS & 5 & 0 & 1.5 & 1 & 0.5$^\ddagger$ & 0.5$^\ddagger$ & \textbf{$\sim$8.5} \\
\hline
\rowcolor{rowalt}
\textbf{Option 1 + 4 combined (best case)} & \textbf{5} & \textbf{0} & \textbf{1} & \textbf{1} & \textbf{0.5} & \textbf{0.5} & \textbf{$\sim$8} \\
\hline
\end{tabularx}
\caption{Latency comparison (5-second utterance). $^*$LLM saved $\sim$0 net due to slow local inference. $^\dagger$Lower quality responses. $^\ddagger$Perceived latency to \emph{first} audio heard, not full response.}
\end{table}

% ════════════════════════════════════════════════════════════════════════════
\section{Honest Assessment of the 1--2 Second Target}

\begin{constraintbox}
\textbf{The 1--2 second total latency target is not achievable with the current architecture.}

Even with every optimisation applied simultaneously:
\begin{itemize}[leftmargin=1.5em]
  \item A 5-second utterance cannot be processed in under 5 seconds — the user must finish
        speaking before we have audio to transcribe
  \item Any cloud STT call adds $\sim$0.5--1.5\,s of network round-trip time
  \item Any cloud LLM call adds $\sim$1--3\,s
  \item Any cloud TTS call adds $\sim$0.5--2\,s
  \item BLE audio playback back to the device adds $\sim$0.5--2\,s
\end{itemize}

The minimum achievable end-to-end latency (time from button release to first audio heard) with
the current ESP32 + BLE + cloud AI architecture is approximately \textbf{3--5 seconds after
speech ends}. For a 5-second utterance this means \textbf{8--10 seconds total}.

This is the standard latency class for consumer voice assistants (Alexa, Google Assistant) and
is considered acceptable for a v1 product. The 1--2 second target is only achievable with
dedicated DSP hardware, always-on wake-word detection, on-device inference chips, and
datacenter-proximity cloud infrastructure — none of which apply here.
\end{constraintbox}

% ════════════════════════════════════════════════════════════════════════════
\section{Recommendation \& Implementation Priority}

\subsection{Phase 1 — Immediate (before PCB finalisation)}

\begin{enumerate}
  \item \textbf{Option 1: Firmware BLE streaming during recording}
  \begin{itemize}
    \item Refactor the ESP32 recording loop to stream I2S chunks directly over BLE notifications
          instead of writing to LittleFS flash
    \item Remove the \texttt{NEW:} / \texttt{SEND:} / \texttt{FILE:} protocol; replace with
          \texttt{REC\_START} / raw binary chunks / \texttt{REC\_END} markers
    \item App-side: accumulate chunks, run STT pipeline on \texttt{REC\_END}
    \item \textbf{Expected gain: eliminates 3--5\,s BLE transfer delay}
  \end{itemize}

  \item \textbf{Option 3: MCU-local processing (parallel evaluation)}
  \begin{itemize}
    \item Begin evaluating a companion edge AI module (e.g.\ Raspberry Pi, ESP32-S3 with
          dedicated AI accelerator) running quantised STT, LLM, and TTS locally
    \item This must be decided \emph{before} PCB finalisation — it requires hardware changes
    \item Run in parallel with Option 1 so both approaches can be compared on current hardware
    \item \textbf{Expected outcome: fully offline pipeline with no cloud dependency}
  \end{itemize}
\end{enumerate}

\subsection{Phase 2 — After Phase 1 is stable and tested}

\begin{enumerate}
  \item \textbf{Option 4: Sentence-level overlapped TTS}
  \begin{itemize}
    \item Enable Gemini token streaming via SSE in the LLM edge function
    \item App detects first sentence boundary and triggers TTS immediately
    \item ESP32 begins playing audio while LLM is still generating remaining sentences
    \item \textbf{Expected gain: perceived first-audio latency drops to $\sim$1.5--3\,s
          after button release}
  \end{itemize}
\end{enumerate}

\subsection{Phase 3 — Future consideration (v2)}

\begin{enumerate}
  \item \textbf{On-device Whisper.cpp} — if privacy or STT cost becomes a requirement,
        integrate \texttt{whisper.rn} (the stub is already in the codebase). Whisper Tiny
        ($\sim$40\,MB) achieves $\sim$0.3--0.8\,s on modern phones.
\end{enumerate}

\subsection{Options Not Recommended}

\begin{itemize}
  \item \textbf{On-device LLM (Option 2):} App size (1--4\,GB) is unacceptable for consumer
        deployment. Net latency improvement is negligible due to slow on-device inference.
  \item \textbf{Streaming STT via Supabase:} Supabase edge functions are stateless HTTP —
        they cannot maintain the WebSocket or gRPC connection required by Google STT Streaming
        or Deepgram. This would require adding a separate dedicated server (significant
        infrastructure cost and complexity).
  \item \textbf{Batch partial STT during recording:} Google STT batch API returns no
        transcription results for audio shorter than $\sim$0.5\,s. Sending sub-second chunks
        while recording produces empty responses and does not reduce latency.
\end{itemize}

% ════════════════════════════════════════════════════════════════════════════
\section{Summary Table}

\begin{table}[h!]
\centering
\renewcommand{\arraystretch}{1.5}
\begin{tabularx}{\linewidth}{|l|X|c|c|c|}
\hline
\rowcolor{primary}
\textcolor{white}{\textbf{\#}} &
\textcolor{white}{\textbf{Latency gain}} &
\textcolor{white}{\textbf{HW change}} &
\textcolor{white}{\textbf{App change}} &
\textcolor{white}{\textbf{Priority}} \\
\hline
Opt 1 & High ($-$3--5\,s post-speech) & Yes & Minor & \textbf{Phase 1} \\
\hline
\rowcolor{rowalt}
Opt 3 & Low — quality trade-off; fully offline & Major & Minor & \textbf{Phase 1} \\
\hline
Opt 4 & High (perceived, first audio in 1.5--3\,s) & No & Major & \textbf{Phase 2} \\
\hline
\rowcolor{rowalt}
Whisper & Medium ($-$1--1.5\,s STT) & No & Major & Phase 3 \\
\hline
Opt 2 & Very Low (net) & No & Major & \textbf{Not recommended} \\
\hline
\rowcolor{rowalt}
Stream STT & Medium ($-$1\,s) needs new server & No & Major & \textbf{Not recommended v1} \\
\hline
\end{tabularx}
\caption{Option summary and implementation priority}
\end{table}

% ════════════════════════════════════════════════════════════════════════════
% ════════════════════════════════════════════════════════════════════════════
\section{Client Follow-Up — Gemini Live Assessment}

The following response was received after the initial report was shared:

\begin{clientbox}
\textbf{From:} Kiran \hfill \textbf{To:} Development Team\\[6pt]

Thank you for the detailed report from the team, appreciate the effort to reduce the latency.

I want to see if we can reduce latency even further below 3--5\,s. I'd like us to evaluate
\textbf{Gemini Live} as a possible production-ready realtime voice architecture as I believe we
can get into the 1--2\,s range if we use it instead potentially.

My understanding is that Gemini Live is designed for continuous low-latency audio sessions over
a stateful WebSocket, where audio is sent in progressive chunks and audio responses are also
returned in chunks. The intended flow would be:

\begin{itemize}[leftmargin=1.5em]
  \item Toy captures audio
  \item Toy sends progressive BLE audio chunks to the phone
  \item Phone forwards those chunks in real time to Gemini Live
  \item Gemini Live returns transcripts / audio response chunks
  \item Phone either plays them directly or sends them back to the toy for speaker playback
\end{itemize}

So the BLE progressive chunking approach still remains important — the difference is that the
phone would then forward those chunks into a realtime Gemini Live session instead of using a
slower request/response pipeline.

Please assess whether this is technically feasible with our current toy + phone setup, what
changes would be required, and whether this would materially improve latency.

\medskip
\textbf{Kiran}
\end{clientbox}

% ─────────────────────────────────────────────────────────────────────────────
\subsection{What is Gemini Live?}

Gemini Live (also referred to as the Gemini Multimodal Live API) is Google's real-time,
bidirectional audio API built on top of Gemini 2.0 Flash. Unlike the separate
STT~$\rightarrow$~LLM~$\rightarrow$~TTS pipeline currently in use, Gemini Live accepts
\textbf{raw audio in} and returns \textbf{raw audio out} over a single persistent WebSocket
connection. There is no separate transcription step, no separate LLM request, and no separate
TTS call — the entire voice conversation happens inside one stateful session.

Key technical properties relevant to this project:

\begin{table}[h!]
\centering
\renewcommand{\arraystretch}{1.3}
\begin{tabularx}{\linewidth}{|l|X|}
\hline
\rowcolor{primary}
\textcolor{white}{\textbf{Property}} & \textcolor{white}{\textbf{Value}} \\
\hline
Transport               & Persistent WebSocket (bidirectional) \\
\hline
\rowcolor{rowalt}
Audio input format      & PCM 16-bit, 16\,kHz, mono — \textbf{identical to ESP32 I2S output} \\
\hline
Audio output format     & PCM 16-bit, 24\,kHz, mono \\
\hline
\rowcolor{rowalt}
Response start latency  & $\sim$300--800\,ms after end of speech detected \\
\hline
Built-in VAD            & Yes — model detects end of utterance automatically \\
\hline
\rowcolor{rowalt}
Conversation memory     & Yes — full multi-turn session state maintained \\
\hline
Underlying model        & Gemini 2.0 Flash (same intelligence as current Gemini Pro call) \\
\hline
\rowcolor{rowalt}
Barge-in support        & Yes — user can interrupt the AI mid-response \\
\hline
\end{tabularx}
\caption{Gemini Live technical properties}
\end{table}

% ─────────────────────────────────────────────────────────────────────────────
\subsection{Feasibility Assessment}

\textbf{Short answer: Yes — Gemini Live is technically feasible and is the recommended path
to achieve the 1--2 second post-speech latency target.}

The audio format is a perfect match. The ESP32 currently records 16\,kHz, 16-bit, mono PCM —
exactly the format Gemini Live expects as input. No re-encoding or sample-rate conversion is
needed on the firmware or the app side for the audio going \emph{in}. The response audio
(24\,kHz PCM) may require a simple downsampling step before playback through the ESP32 speaker
DAC, which is a minor app-side operation.

The critical architectural change is that the current request--response pipeline (three separate
cloud edge function calls: STT, LLM, TTS) is \textbf{replaced entirely} by a single persistent
WebSocket session. This eliminates three sequential network round trips and replaces them with
one continuous stream.

% ─────────────────────────────────────────────────────────────────────────────
\subsection{Required Changes}

\subsubsection{Firmware (ESP32)}

\begin{itemize}
  \item \textbf{Option 1 BLE streaming is a prerequisite} — audio must stream over BLE in
        real time rather than being transferred after recording. No additional firmware change
        beyond Option 1 is required for the audio capture side.
  \item The audio output path (receiving and playing the AI response) remains unchanged —
        the app still sends PCM chunks over BLE and the ESP32 plays them through the DAC.
  \item Optional: remove LittleFS dependency entirely once BLE streaming is in place.
\end{itemize}

\subsubsection{Mobile App}

\begin{itemize}
  \item \textbf{New \texttt{GeminiLiveService}:} Manages a persistent WebSocket connection to
        Gemini Live. Opens the session on app launch; reconnects automatically on drop.
  \item \textbf{BLE chunk forwarding:} Each BLE audio chunk received from the ESP32 is
        immediately written to the Gemini Live WebSocket as a \texttt{realtimeInput} message
        — no buffering, no assembly into a full WAV file.
  \item \textbf{Response playback:} Audio chunks returned by Gemini Live are queued and
        streamed to the ESP32 over BLE as they arrive, without waiting for the full response.
  \item \textbf{Remove STT, LLM, and TTS services:} The \texttt{GoogleSTTService},
        \texttt{LLMService}, and \texttt{TTSService} edge function calls are replaced by the
        single Gemini Live session. The Supabase edge functions for these three steps become
        redundant.
  \item \textbf{Personality / system prompt:} The toy's custom personality (currently sent as
        an LLM system prompt) is sent as the Gemini Live session setup instruction at session
        open — no change to the product behaviour.
\end{itemize}

\subsubsection{Backend Infrastructure}

This is the one area where Gemini Live introduces a constraint that the current Supabase-only
architecture cannot directly satisfy.

\begin{constraintbox}
\textbf{Supabase edge functions cannot maintain a persistent WebSocket connection to Gemini
Live.} Edge functions are stateless HTTP handlers — they terminate after each request. A
Gemini Live session is a long-lived WebSocket that must remain open for the duration of the
conversation (potentially minutes). This is architecturally incompatible with Supabase edge
functions.
\end{constraintbox}

Two approaches are available:

\begin{enumerate}
  \item \textbf{Direct connection from app (v1 / prototype):} The phone app connects to the
        Gemini Live WebSocket directly, with the Gemini API key embedded in the app or fetched
        from Supabase at session start. This is acceptable for a development prototype and
        internal testing. It carries an API key exposure risk and should not be used for
        public release without key rotation controls.

  \item \textbf{Lightweight WebSocket proxy server (production):} A small Node.js or Python
        server (deployable to Google Cloud Run, Railway, Fly.io, or Render — all have free
        tiers) acts as a relay: it accepts the WebSocket from the phone, holds the Gemini API
        key securely, and proxies audio frames in both directions to Gemini Live. This server
        is stateless in terms of data (it just relays bytes) and can be deployed in minutes.
        Supabase continues to handle authentication, database, and user management — only the
        real-time audio relay moves to the proxy.
\end{enumerate}

% ─────────────────────────────────────────────────────────────────────────────
\subsection{Latency Impact}

\begin{table}[h!]
\centering
\renewcommand{\arraystretch}{1.4}
\begin{tabularx}{\linewidth}{|l|X|r|r|}
\hline
\rowcolor{primary}
\textcolor{white}{\textbf{Stage}} &
\textcolor{white}{\textbf{Description}} &
\textcolor{white}{\textbf{Current}} &
\textcolor{white}{\textbf{Gemini Live}} \\
\hline
Recording + BLE stream  & User speaks; chunks sent in real time  & 5.0\,s & 5.0\,s \\
\hline
\rowcolor{rowalt}
Post-recording transfer & BLE file transfer after recording ends & $\sim$3--5\,s & \textbf{0\,s} \\
\hline
STT                     & Speech-to-text cloud call              & $\sim$1--2\,s & \textbf{0\,s} \\
\hline
\rowcolor{rowalt}
LLM                     & Language model inference               & $\sim$1--3\,s & \textbf{0\,s} \\
\hline
TTS                     & Text-to-speech synthesis               & $\sim$1--3\,s & \textbf{0\,s} \\
\hline
\rowcolor{rowalt}
Gemini Live response    & Integrated STT+LLM+TTS — first audio chunk & n/a & $\sim$0.3--0.8\,s \\
\hline
BLE audio to toy        & Audio chunks streamed to ESP32         & $\sim$1--3\,s & $\sim$0.5--1.5\,s \\
\hline
\rowcolor{primary}
\textcolor{white}{\textbf{Post-speech latency}} &
\textcolor{white}{\textbf{Time from button release to first audio heard}} &
\textcolor{white}{\textbf{$\sim$7--16\,s}} &
\textcolor{white}{\textbf{$\sim$0.8--2.3\,s}} \\
\hline
\end{tabularx}
\caption{Latency comparison: current pipeline vs Gemini Live (5-second utterance)}
\end{table}

The post-speech latency with Gemini Live is estimated at \textbf{0.8--2.3 seconds} — directly
within the 1--2 second target range the client has requested.

% ─────────────────────────────────────────────────────────────────────────────
\subsection{Pros and Cons of Gemini Live}

\begin{table}[h!]
\centering
\renewcommand{\arraystretch}{1.3}
\begin{tabularx}{\linewidth}{|l|X|}
\hline
\rowcolor{primary}
\textcolor{white}{\textbf{Pros}} & \textcolor{white}{\textbf{Detail}} \\
\hline
Latency & Post-speech response in $\sim$0.8--2.3\,s — meets the 1--2\,s target \\
\hline
\rowcolor{rowalt}
Architecture simplicity & Replaces three separate cloud services (STT, LLM, TTS) with one WebSocket \\
\hline
Audio format match & ESP32 I2S output (16\,kHz PCM) is the exact required input format \\
\hline
\rowcolor{rowalt}
Built-in VAD & Model detects end of speech — removes need for hardware VAD or button logic \\
\hline
Barge-in & User can interrupt the toy mid-sentence naturally \\
\hline
\rowcolor{rowalt}
Conversation memory & Multi-turn context maintained within a session automatically \\
\hline
\end{tabularx}

\vspace{6pt}

\begin{tabularx}{\linewidth}{|l|X|}
\hline
\rowcolor{warning}
\textcolor{white}{\textbf{Cons / Risks}} & \textcolor{white}{\textbf{Detail}} \\
\hline
Proxy server required & Supabase edge functions cannot maintain a persistent WebSocket; a lightweight relay server is needed for production \\
\hline
\rowcolor{rowalt}
Custom voice lost & Resemble AI custom voice is replaced by Google's built-in TTS voices — personality can still be set via system prompt \\
\hline
API availability & Gemini Live is available on Gemini 2.0 Flash; API stability and pricing should be confirmed before committing \\
\hline
\rowcolor{rowalt}
Audio output resampling & Response audio is 24\,kHz; minor app-side downsampling needed if the ESP32 DAC is configured for 16\,kHz \\
\hline
Significant app rewrite & GeminiLiveService replaces three existing services; requires careful testing \\
\hline
\end{tabularx}
\caption{Gemini Live pros and cons}
\end{table}

% ─────────────────────────────────────────────────────────────────────────────
\subsection{Recommendation on Gemini Live}

\begin{mdframed}[backgroundcolor=hwbg, linecolor=primary, linewidth=1.5pt, roundcorner=4pt,
  innerleftmargin=12pt, innerrightmargin=12pt, innertopmargin=10pt, innerbottommargin=10pt]
\textbf{Gemini Live is the recommended path if the 1--2 second latency target is a firm
requirement for v1.}

\medskip
The proposed architecture is sound. The BLE progressive chunking (Option 1 firmware) remains
a prerequisite and is unchanged. The key shift is replacing the three-step cloud pipeline with
a single Gemini Live WebSocket session. This is technically feasible on the current hardware
with no firmware changes beyond Option~1.

\medskip
\textbf{Suggested next step:} Implement a direct-connection prototype (phone $\rightarrow$
Gemini Live, API key fetched from Supabase at session open) to validate latency on real
hardware. If latency targets are confirmed, proceed to build the lightweight WebSocket proxy
server for production. Total additional development effort is estimated at 3--5 days for the
prototype and 2--3 days for the production proxy server.
\end{mdframed}

% ════════════════════════════════════════════════════════════════════════════
\section*{Conclusion}

The most impactful single change is the \textbf{firmware-side BLE streaming} during recording
(Option~1), which eliminates the largest source of latency — the post-recording file transfer.
Combined with sentence-level overlapped TTS (Option~4) as a Phase~2 addition, the user will
perceive a response within \textbf{1.5--3 seconds of finishing speaking}.

The total end-to-end time from button press to first audio heard is approximately
\textbf{8--10 seconds} — however, it is important to note that \textbf{5 of those seconds are
the recording itself} (the time the user is actively speaking). The actual processing delay
\emph{after} the user finishes speaking is only \textbf{3--5 seconds}, which is comparable to
mainstream consumer voice assistants such as Alexa and Google Assistant. The client should
evaluate the acceptance criterion as \textbf{post-speech latency} (time from button release to
first audio heard), not total elapsed time from button press.

\medskip
Following the client's follow-up request, \textbf{Gemini Live is assessed as technically
feasible and is the clearest path to the 1--2 second post-speech latency target.} The ESP32
audio format is a direct match for the Gemini Live input requirements. A prototype can be
built and tested on current hardware within days. If the prototype validates the latency
estimates, Gemini Live should be adopted as the v1 production voice architecture, with a
lightweight WebSocket proxy server added for API key security before public release.

\end{document}
